\input{../tex-inputs/latex-leseansicht-vorspann}

\section[Arthur Schnitzler an Gustav Schwarzkopf, 19. 7. 1899]{L04420 Arthur Schnitzler an Gustav Schwarzkopf, 19. 7. 1899}
\nopagebreak\mylabel{L04420v}
\rehead{ }\normalsize\beginnumbering\briefempfaengerindex{Schwarzkopf, Gustav@\textsc{Schwarzkopf, Gustav}!zzzSchnitzler, Arthur@\emph{von Arthur Schnitzler}!1899-07-191@{19. 7. 1899}|(be}
\toendnotes[C]{\smallbreak\pagebreak[2]}
\correspDesc{Versand  durch Arthur Schnitzler am 19. 7. 1899 in Velden am Wörthersee
\newline{}Übermittlung  am 19. 7. 1899 in Velden am Wörthersee
\newline{}Zustellung  am 20. 7. 1899 in Wien
\newline{}Erhalt  durch Gustav Schwarzkopf im Zeitraum [20. 7. 1899
                  – 24. 7. 1899?] in Wien}\toendnotes[C]{\smallbreak}
\Standort{DLA, A:Schnitzler, HS.1985.1.1897.}
\physDesc{Bildpostkarte, 266 Zeichen
\newline{}Handschrift: Bleistift, deutsche Kurrent
\newline{}Versand: 1) Stempel: »\nobreak{}\oindex{Velden am Wörthersee@\textbf{Velden am Wörthersee}|pwk}Velden am Wörthersee, 19/7 9\textcolor{gray}{9}, \textcolor{gray}{8}N\nobreak{}«.   2) Stempel: »\nobreak{}\oindex{I., Innere Stadt@\textbf{I., Innere Stadt}|pwk}Wien 1/1 1, 20. 7. 99, 10–11½V., Bestellt\nobreak{}«. }\toendnotes[C]{\smallbreak}\pstart{}{\pb}Herrn \textsc{Gustav Schwarzkopf}\pend{}\pstart{}Wien\oindex{Wien@\textbf{Wien}|pw}\pend{}\pstart{}\textsc{\textsc{I. Tiefer Graben 23}}\oindex{Wien@\textbf{Wien}!I., Innere Stadt@\textbf{I., Innere Stadt}!Tiefer Graben 23@\textbf{Tiefer Graben 23}|pw}.\pend{}{\bigskip}
\pstart
           \centering{}{\pb}\textcolor{gray}{\textbf{Gruss aus Velden\oindex{Velden am Wörthersee@\textbf{Velden am Wörthersee}|pw}.}}\pend
           \vspace{1em}
\pstart
           \raggedleft{}{\pb}Mittwoch Abend.\pend
           \vspace{0.5em}
\pstart
           Adreſſe bleibt bis auf weiteres \textsc{Pension Pundschu\oindex{Pension Pundschu@\textbf{Pension Pundschu}|pw}}.\pend
           
\pstart
           – \label{K_L04420-1v}\edtext{Geſtern}{\lemma{\textnormal{\emph{Gestern}}}\Cendnote{\textnormal{Vgl. A. S.: \emph{Wiener Schnitzler}, 18. 7. 1899.}}}\label{K_L04420-1} früh hier angeko{\geminationm}en. – \label{K_L04420-2v}\edtext{Geſtern{ }Nachmittg}{\lemma{\textnormal{\emph{Gestern Nachmittg}}}\Cendnote{\textnormal{Vgl. A. S.: \emph{Tagebuch}, 18. 7. 1899.}}}\label{K_L04420-2} und
                  \label{K_L04420-3v}\edtext{heute auf einer{ }ſchönen Radpartie}{\lemma{\textnormal{\emph{heute … Radpartie}}}\Cendnote{\textnormal{Vgl. A. S.: \emph{Tagebuch}, 19. 7. 1899.}}}\label{K_L04420-3} mit Burger’s\pwindex{Burger, Caroline 11.\,7.\,1869 Wien – 15.\,3.\,1959 ebd.@\textsc{Burger, Caroline} (11.\,7.\,1869 Wien – 15.\,3.\,1959 ebd.)|pw}\pwindex{Burger, Rudolf *~6.\,12.\,1866 Wien@\textsc{Burger, Rudolf} (*~6.\,12.\,1866 Wien)|pw}. – \label{K_L04420-4v}\edtext{Richard\pwindex{Beer-Hofmann, Richard 11.\,7.\,1866 Wien – 26.\,9.\,1945 New York City@\textsc{Beer-Hofmann, Richard} (11.\,7.\,1866 Wien – 26.\,9.\,1945 New York City)|pw} arbeitet und wird erſt So{\geminationn}tag}{\lemma{\textnormal{\emph{Richard … Sonntag}}}\Cendnote{\textnormal{Vgl. XXXX Auszeichnungsfehler: Dokument L00946 nicht gefunden.}}}\label{K_L04420-4} ſichtbar.\pend
           \pstart Herzlichſt Ihr \spacefill\mbox{A.}\pend{}\selectlanguage{ngerman}\endnumbering\briefempfaengerindex{Schwarzkopf, Gustav@\textsc{Schwarzkopf, Gustav}!zzzSchnitzler, Arthur@\emph{von Arthur Schnitzler}!1899-07-191@{19. 7. 1899}|)be}\mylabel{L04420h}
\begin{anhang}
\end{anhang}\newcommand{\dateiname}{L04420}\newcommand{\titel}{Arthur Schnitzler an Gustav Schwarzkopf, 19. 7. 1899}\newcommand{\editorInnen}{Herausgegeben von Jahnke, SelmaMüller, Martin Anton}\input{../tex-inputs/latex-leseansicht-abspann}
